% Options for packages loaded elsewhere
\PassOptionsToPackage{unicode}{hyperref}
\PassOptionsToPackage{hyphens}{url}
\documentclass[
  ignorenonframetext,
]{beamer}
\newif\ifbibliography
\usepackage{pgfpages}
\setbeamertemplate{caption}[numbered]
\setbeamertemplate{caption label separator}{: }
\setbeamercolor{caption name}{fg=normal text.fg}
\beamertemplatenavigationsymbolsempty
% remove section numbering
\setbeamertemplate{part page}{
  \centering
  \begin{beamercolorbox}[sep=16pt,center]{part title}
    \usebeamerfont{part title}\insertpart\par
  \end{beamercolorbox}
}
\setbeamertemplate{section page}{
  \centering
  \begin{beamercolorbox}[sep=12pt,center]{section title}
    \usebeamerfont{section title}\insertsection\par
  \end{beamercolorbox}
}
\setbeamertemplate{subsection page}{
  \centering
  \begin{beamercolorbox}[sep=8pt,center]{subsection title}
    \usebeamerfont{subsection title}\insertsubsection\par
  \end{beamercolorbox}
}
% Prevent slide breaks in the middle of a paragraph
\widowpenalties 1 10000
\raggedbottom
\AtBeginPart{
  \frame{\partpage}
}
\AtBeginSection{
  \ifbibliography
  \else
    \frame{\sectionpage}
  \fi
}
\AtBeginSubsection{
  \frame{\subsectionpage}
}
\usepackage{iftex}
\ifPDFTeX
  \usepackage[T1]{fontenc}
  \usepackage[utf8]{inputenc}
  \usepackage{textcomp} % provide euro and other symbols
\else % if luatex or xetex
  \usepackage{unicode-math} % this also loads fontspec
  \defaultfontfeatures{Scale=MatchLowercase}
  \defaultfontfeatures[\rmfamily]{Ligatures=TeX,Scale=1}
\fi
\usepackage{lmodern}
\ifPDFTeX\else
  % xetex/luatex font selection
\fi
% Use upquote if available, for straight quotes in verbatim environments
\IfFileExists{upquote.sty}{\usepackage{upquote}}{}
\IfFileExists{microtype.sty}{% use microtype if available
  \usepackage[]{microtype}
  \UseMicrotypeSet[protrusion]{basicmath} % disable protrusion for tt fonts
}{}
\makeatletter
\@ifundefined{KOMAClassName}{% if non-KOMA class
  \IfFileExists{parskip.sty}{%
    \usepackage{parskip}
  }{% else
    \setlength{\parindent}{0pt}
    \setlength{\parskip}{6pt plus 2pt minus 1pt}}
}{% if KOMA class
  \KOMAoptions{parskip=half}}
\makeatother
\setlength{\emergencystretch}{3em} % prevent overfull lines
\providecommand{\tightlist}{%
  \setlength{\itemsep}{0pt}\setlength{\parskip}{0pt}}
\usepackage{bookmark}
\IfFileExists{xurl.sty}{\usepackage{xurl}}{} % add URL line breaks if available
\urlstyle{same}
\hypersetup{
  hidelinks,
  pdfcreator={LaTeX via pandoc}}

\author{\texorpdfstring{}{}}
\date{}

\begin{document}

\section{Methodology: Program Graphs and Intermediate
Representations}\label{methodology-program-graphs-and-intermediate-representations}

\begin{frame}[fragile]{Program graph}
\protect\phantomsection\label{program-graph}
Let \(G = (V, E)\) denote a directed graph whose vertices \(V\)
represent program entities and whose edges \(E\) represent
relationships. COGANT assigns each node a \textbf{kind} (for example
function, variable, type) and a \textbf{semantic role} (for example
definition versus use). Each node and edge may carry:

\begin{itemize}
\tightlist
\item
  A \textbf{stable identifier} for persistence across runs when hashing
  allows.
\item
  \textbf{Type strings} and attributes recovered from the front end.
\item
  A \textbf{confidence} score in \([0,1]\) and \textbf{provenance}
  metadata (source code, type system, control flow, heuristic, external
  tool).
\end{itemize}

This follows the same philosophical line as code property graphs that
fuse AST, control flow, and data flow into one analyzable structure
{[}@yamaguchi2014modeling{]}, but orients the representation toward
\textbf{Generalized Notation Notation (GNN) export} and, optionally,
tensorized views of the same graph: node kinds and roles map to discrete
indices in both forms, and optional embeddings can augment features as
described in \texttt{../cogant/docs/GNN\_EXPORT.md}.

\begin{block}{Structural equivalence}
\protect\phantomsection\label{structural-equivalence}
Two program graphs \(G_1 = (V_1, E_1)\) and \(G_2 = (V_2, E_2)\) are
usefully compared via typed graph isomorphism: a bijection
\(\phi : V_1 \to V_2\) preserving edge structure and relevant labels
supports deduplication and cross-repository linking when interfaces
align.

\begin{equation}
\label{eq:typed-iso}
(u, v) \in E_1 \iff (\phi(u), \phi(v)) \in E_2
\end{equation}

Equation \ref{eq:typed-iso} formalizes the structural invariant we check
during deduplication and cross-repository linking: a candidate bijection
\(\phi\) is accepted only when every edge in \(G_1\) has a corresponding
edge between the images of its endpoints in \(G_2\) (and vice versa), so
that label-preserving matches strictly preserve the adjacency structure
of both program graphs.
\end{block}
\end{frame}

\begin{frame}[fragile]{Progressive IRs}
\protect\phantomsection\label{progressive-irs}
Processing proceeds through a sequence of representations (see
\texttt{../cogant/docs/SPEC.md}):

\begin{enumerate}
\tightlist
\item
  \textbf{Repo IR} --- entities and relationships extracted from
  parsers.
\item
  \textbf{Program graph IR} --- consolidated graph with deduplication
  and metadata.
\item
  \textbf{Semantic mapping IR} --- output of the translation rule
  engine.
\item
  \textbf{State space IR} --- variables, actions, transitions,
  observations.
\item
  \textbf{Process model IR} --- higher-level control patterns where
  implemented.
\item
  \textbf{Validation IR} --- coverage, confidence analysis, schema
  checks.
\end{enumerate}

Not every stage is equally complete for every repository; the SPEC marks
partial areas (translation rules, state space, Rust acceleration)
explicitly.
\end{frame}

\begin{frame}[fragile]{Translation rules}
\protect\phantomsection\label{translation-rules}
The \textbf{translation} stage applies declarative rules that refine
roles, attach labels, and adjust confidence. Concurrency targets and
layering are described in \texttt{../cogant/docs/ARCHITECTURE.md}. The
rule engine composes passes over the graph in a \textbf{fixpoint loop}:
rules are re-applied until no new semantic mappings emerge or a
configurable iteration cap is reached. Overlapping rules may require
priority ordering as documented there; when multiple rules claim the
same graph fragment, the engine retains the mapping with the highest
confidence score, following the principle that edit representations
should be composable {[}@yin2019learning{]}.

\begin{block}{Fixpoint iteration, conflict resolution, and coverage}
\protect\phantomsection\label{fixpoint-iteration-conflict-resolution-and-coverage}
The shipped \texttt{TranslationEngine} in
\texttt{../cogant/py/cogant/translate/engine.py} realizes the fixpoint
loop concretely. Each iteration walks every registered rule once: the
engine calls \texttt{rule.matches(graph,\ query)} to collect candidate
fragments, then calls \texttt{rule.apply(graph,\ match)} on each,
accumulating any resulting \texttt{SemanticMapping} objects keyed by
their stable IDs. A per-pass counter tracks mappings that were genuinely
new (not already present in the running set), and the loop terminates as
soon as a pass completes with zero additions. The engine logs each
iteration boundary through an internal match log, so the number of
passes required to reach a fixed point is directly observable in the
post-run diagnostics. The default iteration cap is
\texttt{max\_iterations\ =\ 10}; in testing, most repositories converge
well before that bound, and the cap exists primarily as a safety valve
against pathological rule sets that could otherwise oscillate
indefinitely.

After fixpoint termination, the engine invokes
\texttt{\_resolve\_conflicts()} to reconcile mappings whose
\texttt{graph\_fragment\_node\_ids} sets overlap. For each overlapping
pair the engine retains the mapping with the higher
\texttt{confidence\_score} and discards the other, logging a
\texttt{conflict\_resolved} event that records the losing ID, the
winning ID, and the specific overlap set. Rule priority is therefore
expressed indirectly through the confidence model rather than through a
separate ordering table: a higher-priority rule expresses its precedence
by yielding mappings with higher scores, which then win in the conflict
resolution pass. A companion entry point,
\texttt{translate\_with\_confidence()}, runs the standard fixpoint loop,
rescores every surviving mapping through the \texttt{ConfidenceModel},
and then re-resolves conflicts so that any ordering shifts induced by
rescoring are honored.

Translation coverage -- the fraction of graph nodes that received at
least one semantic mapping -- is reported by
\texttt{get\_coverage\_report(graph)}. It returns the total node count,
the number of covered nodes, the number of uncovered nodes, a
\texttt{coverage\_percent} value rounded to two decimal places, and the
sorted list of uncovered node IDs. The uncovered list is intentionally
emitted verbatim so that downstream tooling can target unmapped regions
for manual review or rule extension, rather than burying the gap behind
a single aggregate number {[}@allamanis2018survey{]}.

\begin{algorithm}
\caption{Fixpoint translation engine}
\label{alg:fixpoint}
\KwIn{Program graph $G = (V, E)$; rule set $R$; maximum iterations $K$ (default 10)}
\KwOut{Set of semantic mappings $\mathcal{M}$ keyed by stable ID}
$\mathcal{M} \leftarrow \emptyset$\;
\For{$k \leftarrow 1$ \KwTo $K$}{
    $n_{\text{new}} \leftarrow 0$\;
    \ForEach{rule $r \in R$ sorted by $\text{priority}(r)$ descending}{
        \ForEach{match $m \in r.\text{matches}(G)$}{
            $\mu \leftarrow r.\text{apply}(G, m)$\;
            \If{$\mu \neq \bot$ \textbf{and} $\mu.\text{id} \notin \mathcal{M}$}{
                $\mathcal{M} \leftarrow \mathcal{M} \cup \{\mu\}$\;
                $n_{\text{new}} \leftarrow n_{\text{new}} + 1$\;
            }
        }
    }
    \If{$n_{\text{new}} = 0$}{
        \textbf{break} \tcp*{fixed point reached}
    }
}
$\mathcal{M} \leftarrow \textsc{ResolveConflicts}(\mathcal{M})$\;
\Return $\mathcal{M}$\;
\end{algorithm}

Algorithm \ref{alg:fixpoint} summarizes the engine. Termination is
guaranteed because each iteration either produces at least one new
mapping (whose stable ID is then fixed in \(\mathcal{M}\)) or terminates
by the break condition; the outer \(K\) bound serves as a safety valve
for pathological rule sets.

\begin{algorithm}
\caption{Priority-ordered conflict resolution}
\label{alg:conflict}
\KwIn{Mapping set $\mathcal{M}$; priority function $p(\cdot)$}
\KwOut{Reduced mapping set with no overlapping fragments}
Build inverted index $\mathcal{I}: V \to 2^{\mathcal{M}}$ from node IDs to mappings touching them\;
$\mathcal{C} \leftarrow \emptyset$ \tcp*{conflict pairs}
\ForEach{node $v$ with $|\mathcal{I}(v)| \geq 2$}{
    \ForEach{$(\mu_a, \mu_b) \in \binom{\mathcal{I}(v)}{2}$}{
        $\mathcal{C} \leftarrow \mathcal{C} \cup \{(\mu_a, \mu_b)\}$\;
    }
}
$\mathcal{R} \leftarrow \emptyset$ \tcp*{removal set}
\ForEach{$(\mu_a, \mu_b) \in \mathcal{C}$}{
    \If{$\mu_a \in \mathcal{R}$ \textbf{or} $\mu_b \in \mathcal{R}$}{\textbf{continue}}
    $k_a \leftarrow (p(\mu_a), c(\mu_a))$; $k_b \leftarrow (p(\mu_b), c(\mu_b))$\;
    \eIf{$k_a \geq k_b$}{$\mathcal{R} \leftarrow \mathcal{R} \cup \{\mu_b\}$}{$\mathcal{R} \leftarrow \mathcal{R} \cup \{\mu_a\}$}
}
\Return $\mathcal{M} \setminus \mathcal{R}$\;
\end{algorithm}

Algorithm \ref{alg:conflict} detects conflicts via an inverted index in
\(O(\sum_v |\mathcal{I}(v)|^2)\) worst-case time, which is substantially
faster than the naive all-pairs scan for graphs where most nodes carry
at most one mapping.
\end{block}
\end{frame}

\begin{frame}[fragile]{Confidence Scoring and Evidence Tiers}
\protect\phantomsection\label{confidence-scoring-and-evidence-tiers}
The shipped \texttt{ConfidenceModel} in
\texttt{../cogant/py/cogant/translate/confidence.py} computes a scalar
score in \([0,1]\) from:

\begin{itemize}
\tightlist
\item
  Average confidence over provenance records.
\item
  An evidence-diversity term (scaled, capped).
\item
  A parser certainty factor applied multiplicatively.
\item
  Conflict penalties subtracted after scaling.
\end{itemize}

\begin{equation}
\label{eq:confidence-core}
c = \max\left(0,\ \min\left(1,\ (\bar{e} + \delta_d)\cdot \kappa - \pi\right)\right)
\end{equation}

Here \(\bar{e}\) is the mean evidence confidence, \(\delta_d\) is the
diversity bonus (bounded), \(\kappa\) is parser certainty, and \(\pi\)
aggregates conflict penalties. \textbf{Tiers} (for example static-only
versus static-plus-runtime) are assigned from score thresholds and
evidence source tags (\texttt{determine\_confidence\_tier}); see the
same module for named thresholds and enum values. The manuscript does
not duplicate those literals so they cannot drift from code.
\end{frame}

\begin{frame}[fragile]{State space and behavior}
\protect\phantomsection\label{state-space-and-behavior}
Where traces or coverage are available, \textbf{dynamic} extraction
feeds the state-space compiler. The goal is a compact behavioral model:
states, actions, transitions, and observations that sit alongside the
static graph for tasks that require execution-sensitive features.

The shipped \texttt{StateSpaceCompiler} in
\texttt{../cogant/py/cogant/statespace/compiler.py} constructs this
model in several coordinated passes driven by the semantic mappings and
the underlying program graph.

\textbf{State variables} are identified by a
\texttt{StateVariableExtractor} that traverses graph nodes carrying
READS and WRITES edges: any node that participates in a write is a
candidate hidden-state variable, and its type, cardinality, and initial
confidence are derived from the node's recovered type string and the
provenance of the rules that produced the associated mapping.

\textbf{Actions} are extracted from semantic mappings of kind
\texttt{ACTION}, with parameters read from node metadata (typically the
parser-recovered function signature), effects traced through outgoing
WRITES edges on the controller node, and preconditions derived from
parameter lists, docstring directives, and explicit metadata entries.
For method-kind actions, the compiler also walks CONTAINS edges back to
the enclosing class so that instance-level state mutations are
attributed to the correct controller, mirroring how code property graphs
fuse structural and data-flow views {[}@yamaguchi2014modeling{]}.

\textbf{Transitions} are inferred by a cross-reference pass
(\texttt{\_cross\_reference\_actions\_and\_variables}) that, for each
action, partitions its adjacent variables into reads and writes based on
edge kind (WRITES and MUTATES yielding writes; READS and OBSERVES
yielding reads). The resulting \texttt{Transition} object records a
\texttt{source\_state} in which every touched variable is marked
\texttt{"pre"} and a \texttt{target\_state} in which written variables
advance to \texttt{"post"} while read-only variables remain
\texttt{"pre"}; this simple pre/post convention keeps the model faithful
to the static evidence without overcommitting to symbolic value domains
that cannot be recovered from AST analysis alone. Trigger attribution
follows incoming TRIGGERS and CALLS edges so that orchestration flow is
preserved in the behavioral model.

\textbf{Observation modalities} are built from semantic mappings of kind
\texttt{OBSERVATION}: each associated node becomes an
\texttt{ObservationModality} whose modality type (\texttt{log},
\texttt{metric}, \texttt{event}, \texttt{sensor}, or a generic fallback)
is inferred from the mapping's description and the node's name, giving
the GNN export a typed observation channel aligned with the OBSERVES
edges in the program graph.

The \textbf{temporal regime} of the model is determined by a companion
\texttt{TemporalAnalyzer}. It classifies nodes as asynchronous when
their metadata carries \texttt{is\_async}/\texttt{async} flags or their
names match patterns such as \texttt{async}, \texttt{callback},
\texttt{promise}, or \texttt{future}, and as event-related when their
kind is \texttt{EVENT} or their names match \texttt{event},
\texttt{handler}, \texttt{listener}, or \texttt{trigger}. Temporal
orderings are then extracted from CALLS and TRIGGERS edges, with each
edge classified as \texttt{parallel} (when either endpoint is async) or
\texttt{sequential} otherwise, and event patterns are assembled from
triggers-in / triggers-out pairs around each event node. A final
decision rule selects among \texttt{SYNCHRONOUS}, \texttt{ASYNCHRONOUS},
\texttt{EVENT\_DRIVEN}, and \texttt{HYBRID}: the presence of event
triggers and patterns combined with async handlers yields
\texttt{HYBRID}; event triggers alone yield \texttt{EVENT\_DRIVEN}; an
async fraction above 30 percent or the presence of async handlers yields
\texttt{ASYNCHRONOUS}; and all remaining cases default to
\texttt{SYNCHRONOUS}. This regime is attached to the
\texttt{StateSpaceModel} metadata so downstream consumers know which
execution model the transition graph assumes.

When coverage and traces are available, \texttt{enrich\_graph()} in
\texttt{../cogant/py/cogant/dynamic/enrichment.py} feeds additional
evidence into this pipeline. Coverage enrichment matches
\texttt{.coverage} SQLite databases or Cobertura XML reports against
nodes whose \texttt{path} and \texttt{source\_range} overlap covered
lines, attaching \texttt{coverage\_hits} and, where branch data is
available, \texttt{branch\_coverage} metadata. Trace enrichment parses
Chrome DevTools traces, writes \texttt{call\_count},
\texttt{avg\_duration\_ms}, and \texttt{is\_hot\_path} onto matching
callable nodes, and adds or reweights dynamic CALLS edges tagged with
\texttt{evidence\_sources={[}"dynamic\_trace"{]}}. Both steps also
append \texttt{dynamic\_coverage} and \texttt{dynamic\_trace} markers to
the program graph's evidence sources. The confidence model consumes
these markers directly: any mapping whose evidence set now contains both
static and dynamic entries becomes eligible for promotion from the
\texttt{STATIC\_ONLY} tier to \texttt{STATIC\_PLUS\_RUNTIME}, and
hot-path and branch-coverage metadata raise the underlying score through
the diversity bonus \(\delta_d\) in Equation \ref{eq:confidence-core}.
This is the mechanism by which executing the target program, even
partially, converts static heuristics into corroborated behavioral facts
without rerunning the upstream rule engine {[}@allamanis2018survey{]}.
\end{frame}

\begin{frame}[fragile]{GNN export (Generalized Notation Notation)}
\protect\phantomsection\label{gnn-export-generalized-notation-notation}
Exports package the program graph and compiled state-space model into
the Active Inference Institute's \textbf{Generalized Notation Notation}
plus a small set of interop targets:

\begin{itemize}
\tightlist
\item
  \textbf{GNN Markdown (\texttt{model.gnn.md})} --- canonical,
  human-readable Generalized Notation Notation organized into the
  section order defined in \texttt{../cogant/docs/GNN\_EXPORT.md} (Model
  Metadata, Repository Metadata, Source Coverage, State Space,
  Observation Modalities, Actions/Policies, Connections, Factors,
  Transition Structure, Likelihood Structure, Preferences/Constraints,
  Time Settings, Parameterization, Ontology Mapping, Provenance,
  Confidence Scores, Rendering Hints, Validation Notes).
\item
  \textbf{GNN JSON (\texttt{model.gnn.json})} --- machine-readable
  equivalent of the markdown sections, plus companion JSON files per
  section (\texttt{state\_space.json}, \texttt{observations.json},
  \texttt{actions.json}, \texttt{transitions.json}, and so on).
\item
  \textbf{Interop exports} --- GraphML and Parquet for analysis in
  Gephi/yEd and DuckDB, and optional tensor views (PyTorch Geometric
  \texttt{Data} objects, DGL graphs, HDF5 tables) for downstream graph
  neural network training pipelines that consume the program graph as a
  relational tensor.
\end{itemize}

Node and edge feature breakdowns, section contracts, and the optional
framework targets are specified in
\texttt{../cogant/docs/GNN\_EXPORT.md}; they determine the structure of
the emitted notation as well as the effective input dimensionality of
any downstream model.
\end{frame}

\begin{frame}{Error handling philosophy}
\protect\phantomsection\label{error-handling-philosophy}
The implementation distinguishes fatal configuration or parse failures
from per-file errors that allow partial results. Validation aggregates
warnings and inconsistencies into a report that travels with the bundle.
This mirrors the layered error categories documented in the architecture
(fatal, error, warning, info).
\end{frame}

\end{document}
